\section{Every bounded predictable input has the same mean exponent}
\label{bcl:section}

Fix $Q<\infty$, and let $\mathcal Q_Q$ be the nonanticipating policies
that add a leaf at the retained root with predictable intensity
$0\le q_t\le Q$.  Apart from these arrivals the retained component
has the inert generator.  Policies may depend on the observed component
and additional nonanticipating randomness.  They do not anticipate
the future Poisson driving measures.  The component starts from the
inert newborn state $x_\infty$.  Put
\[
 F_q(t)=\mathbb E_q C(t),\qquad
 F_Q^-(t)=\inf_{q\in\mathcal Q_Q}F_q(t),\qquad
 F_Q^+(t)=\sup_{q\in\mathcal Q_Q}F_q(t).
\]
The resolvent estimate controls $F_Q^+$; a preparation and isolation argument gives the matching lower rate uniformly over the same class. No ordering of the tree laws under different input rates is required.

\begin{lemma}[Uniform preparation and isolation of an inert descendant]
\label{bcl:preparation}
For every integer $M\ge2$, every $q\in\mathcal Q_Q$, and every $t\ge0$,
\begin{equation}\label{bcl:uniform-lower}
 \boxed{
 F_q(t+2)\ge p_{M,Q}e^{-t/M}F_\infty(t),\qquad
 p_{M,Q}=\frac{M-1}{M}\,e^{-2Q-5}
                         (1-e^{-1})^{M-2}>0.
 }
\end{equation}
Here $F_\infty$ is the inert newborn mean, already known to have
exponent $\mu$.  Only a fixed preparation interval suppresses the
external input.  No suppression of that input is imposed after time
two.
\end{lemma}
\begin{proof}
Write $u$ for the original retained root.  Realize the controlled
arrivals by thinning an independent rate-$Q$ candidate process,
accepting a candidate with conditional probability $q_t/Q$ when
$Q>0$.  For $Q=0$ this channel is absent.  Suppress all candidates
on $[0,2]$, at probability cost $e^{-2Q}$, independently of the
internal driving clocks.  On this event the preparation is the same
for every admissible policy.

In the first unit interval keep $u$'s clock silent and require its
initial one-leaf Yule family to have exactly $M-1$ leaves at time one.
The root silence has probability $e^{-1}$, and the Yule probability is
\[
 e^{-1}(1-e^{-1})^{M-2}.
\]
The full degree of $u$ at time one is then $M$, counting its inert
stub.  There are no other cores in this prepared component.

In the second unit interval require exactly one ring of $u$'s clock.
This has probability $e^{-1}$.  Before that ring all its real
neighbours are leaves, whose number can only increase, so the
probability of choosing a leaf is at least $(M-1)/M$.  This promotion
creates a core child $v$ with one leaf.  Require $v$ and its new leaf
to have no rings before time two.  Given the promotion time, this has
probability at least $e^{-2}$.  Thus at time two $v$ is exactly a
newborn retained core, attached to the real parent $u$ of degree at
least $M$.  Multiplication of these conditional probability bounds
gives the displayed $p_{M,Q}$.  Other leaves at $u$ are free to grow
throughout the second interval; this cannot decrease its full degree.

Condition on the whole preparation history at time two.  Split $u$'s
future rate-one clock into a rate-$1/M$ candidate stream for the edge
$u\to v$ and a rate-$(1-1/M)$ residual stream, with independent fresh
increments.  At a candidate ring choose $v$ with probability
$M/D_u$.  If the candidate is rejected, choose uniformly among the
other $D_u-1$ neighbours; at a residual ring use the same uniform
choice among those other neighbours.  This is a correct construction
of the original source choices: the rate to $v$ is $1/D_u$, and the
rate to any other neighbour is
\[
 \frac{(1/M)(1-M/D_u)+(1-1/M)}{D_u-1}
 =\frac1{D_u}.
\]
It remains valid forever because degrees are nondecreasing and
$D_u\ge M$.

The probability that this candidate stream has no ring in the next
$t$ units of time is exactly $e^{-t/M}$.  The stream is independent
of the internal drivers of $v$ and its future descendants.  On its
no-ring event there can be no arrival from $u$ into $v$.  The retained
descendant component at $v$ consequently follows exactly the inert
newborn process.  An outward choice $v\to u$ changes only the parent
and is a null transition for this descendant component.  The external
control, even if it reacts to the descendant history, also changes
only $u$ and cannot send an arrival across the edge without a point
of the reserved candidate stream.

Conditional on the preparation history, the expected core reward of
this descendant component multiplied by the no-candidate indicator is
therefore $e^{-t/M}F_\infty(t)$.  Its cores are all counted in the
original retained component.  Averaging over the preparation event
proves~\eqref{bcl:uniform-lower}.
\end{proof}

\begin{theorem}[Matching upper and lower rates uniformly over controls]
\label{bcl:matching}
For every fixed finite $Q\ge0$,
\begin{equation}\label{bcl:inf-sup}
 \boxed{
 \lim_{t\to\infty}\frac{\log F_Q^-(t)}t
 =\lim_{t\to\infty}\frac{\log F_Q^+(t)}t
 =\mu.
 }
\end{equation}
In particular, every single admissible predictable policy has mean
exponent $\mu$.  The conclusion also holds for horizon-dependent
choices of policies with the same common bound $Q$.
\end{theorem}
\begin{proof}
For fixed $M$, Lemma~\ref{bcl:preparation} and the inert exponent give
\[
 \liminf_{t\to\infty}\frac{\log F_Q^-(t)}t
 \ge\mu-\frac1M.
\]
The fixed time shift and the positive preparation probability disappear
at the exponential scale.  Let $M\to\infty$.

For completeness, the upper bound uses the canonical resolvent
$V_s$, not a monotonicity claim.  At any $s>\mu$, write
$K_s=40W(s)/s$ and
\[
 U_s=V_s+\frac{QK_s}{s}.
\]
The established root increment bound and resolvent identity give
\[
 (s-\mathcal L_\infty-q\Delta_{\rm root})U_s\ge C
 \qquad(0\le q\le Q),\qquad U_s\ge C/s.
\]
Stopped Dynkin and nonnegativity give
\[
 F_Q^+(t)\le sU_s(x_\infty)e^{st}.
\]
The controlled population is nonexplosive, since its birth rate is at
most its real population plus $Q$.  Fatou's lemma removes the
localization.  Let $s\downarrow\mu$.  Since
$F_Q^-\le F_q\le F_Q^+$, both the uniform statement and the assertion
for each policy follow.
\end{proof}

\begin{corollary}[A permanently fixed parent has exponent $\mu$]
\label{bcl:fixed-parent}
If the parent has permanently fixed degree $d\ge2$, supplying its
newborn child with independent rate-$1/d$ arrivals, and $F_d(t)$ denotes
the retained mean, then
\[
 \boxed{\lim_{t\to\infty}\frac1t\log F_d(t)=\mu.}
\]
\end{corollary}
\begin{proof}
Take the constant policy $q_t=1/d$ in
Theorem~\ref{bcl:matching}.
\end{proof}

The argument extends to every fixed finite retained starting component
with one inert parent stub and a core root.  Its upper constant is
$sU_s(y)$.  For the lower bound, a finite positive-probability
preparation first creates a root leaf if necessary, then builds a
large root degree and a fresh child as above.  Such a root has at
least one real neighbour, and prescribing a choice of the root by
that neighbour has positive probability.  Suppressing rate-$Q$
candidates during this finite preparation keeps its probability
uniform in the policy.  The resulting constant may depend on the
fixed initial component; the exponent remains $\mu$.

